\begin{tomtat}
	\begin{minipage}[t]{0.3\textwidth}
		\vspace{0pt}
		\begin{tomtat}
			\begin{center}
				\textbf{ĐỀ 33}
			\end{center}
			
		\end{tomtat}
	\end{minipage}
	\hfill
	\begin{minipage}[t]{0.5\textwidth}
		\vspace{0pt}
		\begin{tomtat}
			\begin{center}
				\includegraphics[scale=0.08]{images/33}
			\end{center}
		\end{tomtat}
	\end{minipage}
\end{tomtat}
\textbf{I. PHẦN TRẮC NGHIỆM}
\begin{ex}
	Kết quả của phép tính $A=3\sqrt{2}-\sqrt{8}$ là:
	\choice{$-\sqrt{2}$}{$4\sqrt{2}$}{$\sqrt{2}$}{$6\sqrt{3}$}
	\loigiai{
		Ta có:
		\begin{align*} A &= 3\sqrt{2} - \sqrt{8} \\ &= 3\sqrt{2} - \sqrt{4 \cdot 2} \\ &= 3\sqrt{2} - 2\sqrt{2} \\ &= (3-2)\sqrt{2} \\ &= \sqrt{2} \end{align*}
		\\ \textbf{Đáp Án C}
	}
\end{ex}

\begin{ex}
	Cặp số nào là nghiệm của phương trình: $2x-3y=1$.
	\choice{$(1;1)$}{$(1;-1)$}{$(-1;1)$}{$(-1;-1)$}
	\loigiai{
		Ta thay từng cặp số vào phương trình $2x-3y=1$:
		A. Với $(1;1)$: $2(1) - 3(1) = 2 - 3 = -1 \neq 1$. Loại.
		B. Với $(1;-1)$: $2(1) - 3(-1) = 2 + 3 = 5 \neq 1$. Loại.
		C. Với $(-1;1)$: $2(-1) - 3(1) = -2 - 3 = -5 \neq 1$. Loại.
		D. Với $(-1;-1)$: $2(-1) - 3(-1) = -2 + 3 = 1$. Chọn.
		Vậy cặp số $(-1;-1)$ là nghiệm của phương trình.
		\\ \textbf{Đáp Án D}
	}
\end{ex}

\begin{ex}
	Điều kiện xác định của $\sqrt{x+2}$ là:
	\choice{$x \ge -2$}{$x \ge 2$}{$x \le -2$}{$x \le 2$}
	\loigiai{
		Căn thức $\sqrt{A}$ xác định khi $A \ge 0$.
		Do đó, $\sqrt{x+2}$ xác định khi $x+2 \ge 0$.
		$x+2 \ge 0 \Leftrightarrow x \ge -2$.
		\\ \textbf{Đáp Án A}
	}
\end{ex}

\begin{ex}
	Nghiệm của phương trình: $x^2-5x+4=0$ là:
	\choice{$x=0 \text{ và } x=1$}{$x=1 \text{ và } x=4$}{$x=1 \text{ và } x=3$}{$x=1 \text{ và } x=2$}
	\loigiai{
		Phương trình $x^2-5x+4=0$ có $a=1, b=-5, c=4$.
		Ta thấy $a+b+c = 1 + (-5) + 4 = 0$.
		Do đó, phương trình có hai nghiệm là $x_1 = 1$ và $x_2 = \frac{c}{a} = \frac{4}{1} = 4$.
		Vậy nghiệm của phương trình là $x=1$ và $x=4$.
		\\ \textbf{Đáp Án B}
	}
\end{ex}

\begin{ex}
	Cho tam giác ABC vuông tại A có góc C bằng $40^\circ$. Thì độ dài AB là:
	\choice{$AB=BC \sin 40^\circ$}{$AB=AC \sin 40^\circ$}{$AB=BC \cos 40^\circ$}{$AB=AC \cos 40^\circ$}
	\loigiai{
		Trong tam giác ABC vuông tại A:
		Cạnh AB là cạnh đối của góc C.
		Cạnh BC là cạnh huyền.
		Theo định nghĩa tỉ số lượng giác sin:
		$\sin C = \frac{\text{cạnh đối}}{\text{cạnh huyền}} = \frac{AB}{BC}$.
		$\sin 40^\circ = \frac{AB}{BC}$.
		Suy ra $AB = BC \sin 40^\circ$.
		\\ \textbf{Đáp Án A}
	}
\end{ex}

\begin{ex}
	Hộp sữa có dạng hình trụ có đường kính đáy là 10cm, chiều cao của hộp sữa là 8cm. Thể tích của hộp sữa là:
	\choice{$80 \text{ cm}^3$}{$251,2 \text{ cm}^3$}{$628 \text{ cm}^3$}{$2512 \text{ cm}^3$}
	\loigiai{
		Đường kính đáy của hình trụ là $d = 10$ cm.
		Bán kính đáy của hình trụ là $r = \frac{d}{2} = \frac{10}{2} = 5$ cm.
		Chiều cao của hình trụ là $h = 8$ cm.
		Thể tích của hình trụ được tính bằng công thức $V = \pi r^2 h$.
		$V = \pi \cdot (5)^2 \cdot 8 = \pi \cdot 25 \cdot 8 = 200\pi \text{ (cm}^3\text{)}$.
		Nếu lấy $\pi \approx 3,14$:
		$V \approx 200 \cdot 3,14 = 628 \text{ (cm}^3\text{)}$.
		\\ \textbf{Đáp Án C}
	}
\end{ex}
\begin{ex}
	Trong mặt phẳng tọa độ Oxy cho parabol $(P): y=x^2$ và đường thẳng $(d): y=2x-m+3$. Giá trị m để đường thẳng (d) cắt (P) tại 2 điểm phân biệt là:
	\choice{$m<4$}{$m \ge 4$}{$m \le 4$}{$m>4$}
	\loigiai{
		Phương trình hoành độ giao điểm của (P) và (d) là:
		$x^2 = 2x - m + 3$
		$\Leftrightarrow x^2 - 2x + m - 3 = 0 \quad (*)$
		Để đường thẳng (d) cắt (P) tại 2 điểm phân biệt thì phương trình $(*)$ phải có hai nghiệm phân biệt.
		Điều này xảy ra khi $\Delta' > 0$.
		Ta có $a=1, b'=-1, c=m-3$.
		$\Delta' = (b')^2 - ac = (-1)^2 - 1(m-3)$
		$\Delta' = 1 - m + 3$
		$\Delta' = 4 - m$.
		Để có hai nghiệm phân biệt, $\Delta' > 0$:
		$4 - m > 0$
		$\Leftrightarrow 4 > m$
		$\Leftrightarrow m < 4$.
		\\ \textbf{Đáp Án A}
	}
\end{ex}

\begin{ex}
	Năng suất lúa hè thu (tạ/ha) năm 1998 của 31 tỉnh ở Việt Nam được thống kê trong bảng sau:
	\begin{center}
		\begin{tabular}{|c|c|c|c|c|c|}
			\hline
			Năng suất lúa (Tạ/ha) & 25 & 30 & 35 & 40 & 45 \\
			\hline
			Tần số & 4 & 7 & 9 & 6 & 5 \\
			\hline
		\end{tabular}
	\end{center}
	Giá trị $x_3=35$ có tần số bằng:
	\choice{$6$}{$4$}{$7$}{$9$}
	\loigiai{
		Dựa vào bảng tần số đã cho:
		Khi Năng suất lúa (Tạ/ha) là 25, tần số là 4.
		Khi Năng suất lúa (Tạ/ha) là 30, tần số là 7.
		Khi Năng suất lúa (Tạ/ha) là 35 (đây là $x_3$), tần số là 9.
		Khi Năng suất lúa (Tạ/ha) là 40, tần số là 6.
		Khi Năng suất lúa (Tạ/ha) là 45, tần số là 5.
		Vậy, giá trị $x_3=35$ có tần số bằng 9.
		\\ \textbf{Đáp Án D}
	}
\end{ex}

\newpage
\textbf{II. PHẦN TỰ LUẬN}
\begin{ex}
	Giải hệ phương trình $\begin{cases} 3x+y=5 \\ x-2y=-3 \end{cases}$.
	\loigiai{
		Ta có hệ phương trình:
		$$ \begin{cases} 3x+y=5 \quad &(1) \\ x-2y=-3 \quad &(2) \end{cases} $$
		Từ phương trình (1), suy ra $y = 5-3x$.
		Thế $y = 5-3x$ vào phương trình (2):
		\begin{align*} x - 2(5-3x) &= -3 \\ x - 10 + 6x &= -3 \\ 7x &= -3 + 10 \\ 7x &= 7 \\ x &= 1 \end{align*}
		Thay $x=1$ vào $y = 5-3x$:
		\begin{align*} y &= 5 - 3(1) \\ &= 5 - 3 \\ &= 2 \end{align*}
		Vậy nghiệm của hệ phương trình là $(x;y) = (1;2)$.
	}
\end{ex}

\begin{ex}
	Rút gọn biểu thức sau: $Q=\left(\frac{1}{\sqrt{x}+2}+\frac{\sqrt{x}}{x-4}\right) \cdot \left(\frac{2\sqrt{x}+4}{\sqrt{x}-1}\right)$ (với $x \ge 0, x \neq 1, x \neq 4$).
	\loigiai{
		Điều kiện: $x \ge 0, x \neq 1, x \neq 4$.
		Ta có $x-4 = (\sqrt{x}-2)(\sqrt{x}+2)$.
		\begin{align*} Q &= \left(\frac{1}{\sqrt{x}+2}+\frac{\sqrt{x}}{(\sqrt{x}-2)(\sqrt{x}+2)}\right) \cdot \frac{2(\sqrt{x}+2)}{\sqrt{x}-1} \\ &= \left(\frac{\sqrt{x}-2}{(\sqrt{x}-2)(\sqrt{x}+2)}+\frac{\sqrt{x}}{(\sqrt{x}-2)(\sqrt{x}+2)}\right) \cdot \frac{2(\sqrt{x}+2)}{\sqrt{x}-1} \\ &= \frac{\sqrt{x}-2+\sqrt{x}}{(\sqrt{x}-2)(\sqrt{x}+2)} \cdot \frac{2(\sqrt{x}+2)}{\sqrt{x}-1} \\ &= \frac{2\sqrt{x}-2}{(\sqrt{x}-2)(\sqrt{x}+2)} \cdot \frac{2(\sqrt{x}+2)}{\sqrt{x}-1} \\ &= \frac{2(\sqrt{x}-1)}{(\sqrt{x}-2)(\sqrt{x}+2)} \cdot \frac{2(\sqrt{x}+2)}{\sqrt{x}-1} \end{align*}
		Vì $x \neq 1$ nên $\sqrt{x}-1 \neq 0$. Vì $x \neq 4$ nên $\sqrt{x}+2 \neq 0$.
		\begin{align*} Q &= \frac{2\cancel{(\sqrt{x}-1)}}{(\sqrt{x}-2)\cancel{(\sqrt{x}+2)}} \cdot \frac{2\cancel{(\sqrt{x}+2)}}{\cancel{(\sqrt{x}-1)}} \\ &= \frac{2 \cdot 2}{\sqrt{x}-2} \\ &= \frac{4}{\sqrt{x}-2} \end{align*}
		Vậy $Q = \frac{4}{\sqrt{x}-2}$.
	}
\end{ex}

\begin{ex}
	Cho phương trình: $x^2-x-3=0$, không giải phương trình, hãy tính giá trị biểu thức: $T=(x_1-1)(x_1+1)+(x_2-1)(x_2+1)$.
	\loigiai{
		Phương trình $x^2-x-3=0$ có $a=1, b=-1, c=-3$.
		$\Delta = b^2-4ac = (-1)^2 - 4(1)(-3) = 1+12=13 > 0$.
		Phương trình có hai nghiệm phân biệt $x_1, x_2$.
		Theo hệ thức Vi-ét, ta có:
		$$ \begin{cases} x_1 + x_2 = -\frac{b}{a} = -\frac{-1}{1} = 1 \\ x_1 x_2 = \frac{c}{a} = \frac{-3}{1} = -3 \end{cases} $$
		Ta biến đổi biểu thức T:
		\begin{align*} T &= (x_1-1)(x_1+1)+(x_2-1)(x_2+1) \\ &= (x_1^2 - 1^2) + (x_2^2 - 1^2) \\ &= x_1^2 - 1 + x_2^2 - 1 \\ &= (x_1^2 + x_2^2) - 2 \\ &= ((x_1+x_2)^2 - 2x_1x_2) - 2 \end{align*}
		Thay giá trị từ hệ thức Vi-ét vào biểu thức T:
		\begin{align*} T &= ((1)^2 - 2(-3)) - 2 \\ &= (1 + 6) - 2 \\ &= 7 - 2 \\ &= 5 \end{align*}
		Vậy $T=5$.
	}
\end{ex}
\begin{ex}
	Cho một hộp gồm các thẻ được đánh số từ 1 đến 10, lấy ngẫu nhiên 1 thẻ. Tìm xác suất để thẻ lấy ra được là số chia 3 dư 1.
	\loigiai{
		Không gian mẫu $\Omega = \{1; 2; 3; 4; 5; 6; 7; 8; 9; 10\}$.
		Số phần tử của không gian mẫu là $n(\Omega) = 10$.
		
		Gọi A là biến cố: “Thẻ lấy ra được là số chia 3 dư 1”.
		các kết quả thuận lợi cho biến cố A là $A = \{1; 4; 7; 10\}$.
		Số phần tử của biến cố A là $n(A) = 4$.
		
		Xác suất của biến cố A là:
		$P(A) = \frac{n(A)}{n(\Omega)} = \frac{4}{10} = \frac{2}{5}$.
		\\ \textbf{Đáp số: $P(A) = \frac{2}{5}$}
	}
\end{ex}

\begin{ex}
	Một xí nghiệp sản xuất được 120 sản phẩm loại I và 120 sản phẩm loại II trong thời gian 7 giờ. Mỗi giờ sản xuất được số sản phẩm loại I ít hơn số sản phẩm loại II là 10 sản phẩm. Hỏi mỗi giờ xí nghiệp sản xuất được bao nhiêu sản phẩm mỗi loại?
	\loigiai{
		Gọi số sản phẩm loại I xí nghiệp sản xuất được trong một giờ là $x$ (sản phẩm).
		Gọi số sản phẩm loại II xí nghiệp sản xuất được trong một giờ là $y$ (sản phẩm).
		Điều kiện: $x > 0, y > 0$.
		
		Mỗi giờ sản xuất được số sản phẩm loại I ít hơn số sản phẩm loại II là 10 sản phẩm, nên ta có phương trình:
		$y - x = 10 \implies y = x + 10 \quad (1)$
		
		Thời gian để sản xuất 120 sản phẩm loại I là: $t_1 = \frac{120}{x}$ (giờ).
		Thời gian để sản xuất 120 sản phẩm loại II là: $t_2 = \frac{120}{y}$ (giờ).
		
		Tổng thời gian sản xuất cả hai loại sản phẩm là 7 giờ, nên ta có phương trình:
		$\frac{120}{x} + \frac{120}{y} = 7 \quad (2)$
		
		Thay (1) vào (2):
		\begin{align*} \frac{120}{x} + \frac{120}{x+10} &= 7 \\ 120(x+10) + 120x &= 7x(x+10) \\ 120x + 1200 + 120x &= 7x^2 + 70x \\ 240x + 1200 &= 7x^2 + 70x \\ 7x^2 + 70x - 240x - 1200 &= 0 \\ 7x^2 - 170x - 1200 &= 0 \end{align*}
		Giải phương trình bậc hai $7x^2 - 170x - 1200 = 0$:
		Ta có $\Delta = b^2 - 4ac = (-170)^2 - 4(7)(-1200) = 28900 + 33600 = 62500$.
		$\sqrt{\Delta} = \sqrt{62500} = 250$.
		Phương trình có hai nghiệm phân biệt:
		$x_1 = \frac{-b + \sqrt{\Delta}}{2a} = \frac{170 + 250}{2 \cdot 7} = \frac{420}{14} = 30$.
		$x_2 = \frac{-b - \sqrt{\Delta}}{2a} = \frac{170 - 250}{2 \cdot 7} = \frac{-80}{14} = -\frac{40}{7}$.
		
		Đối chiếu với điều kiện $x > 0$, ta nhận nghiệm $x = 30$.
		Khi đó, $y = x + 10 = 30 + 10 = 40$.
		(Thỏa mãn điều kiện $x > 0, y > 0$).
		
		Vậy, mỗi giờ xí nghiệp sản xuất được 30 sản phẩm loại I và 40 sản phẩm loại II.
		\\ \textbf{Đáp số: Loại I: 30 sản phẩm/giờ, Loại II: 40 sản phẩm/giờ}
	}
\end{ex}
\begin{ex}
	Cho đường tròn tâm O đường kính AB. Vẽ dây cung CD vuông góc với AB tại I (I nằm giữa A và O). Lấy điểm E trên cung nhỏ BC (E khác B và C), AE cắt CD tại F. Chứng minh:
	\begin{enumerate}
		\item[a)] BEFI là tứ giác nội tiếp đường tròn.
		\item[b)] $AE \cdot AF = AC^2$.
	\end{enumerate}
\end{ex}

\begin{ex}
	Cho $x>0, y>0$ và $x+y \ge 6$.
	Tìm giá trị nhỏ nhất của biểu thức: $P = 3x+2y+\frac{6}{x}+\frac{8}{y}$.
\end{ex}

\setcounterref{ex}{0}
